% =============================================================================
% PREFÁCIO — Edição Completa (volume único)
% -----------------------------------------------------------------------------
% Arquivo editável à mão, como a capa. NÃO é gerado por
% scripts/gerar_edicao_completa.py — o gerador apenas o inclui, depois do
% sumário e das listas e antes da página de Parte do Volume 0 (posição
% canônica: ABNT NBR 6029, 4.2.2, onde o prefácio abre os elementos textuais;
% Chicago Manual of Style, 17ª ed., §1.4).
%
% Os prefácios dos volumes avulsos continuam intactos nos seus próprios
% arquivos; na edição reunida eles dão lugar a este.
%
% >>> TEXTO DE PARTIDA, ESCRITO PARA SER SUBSTITUÍDO PELO AUTOR. <<<
% Sem números que envelhecem (contagem de páginas, de equações): o registro do
% projeto mostra que o erro recorrente destes textos não é conteúdo errado, é
% status congelado.
% =============================================================================
\clearpage
\section*{Prefácio à Edição Completa}
\addcontentsline{toc}{section}{Prefácio à Edição Completa}
\markboth{Prefácio à Edição Completa}{}

Esta edição reúne num só volume os três que compõem a \colecaotitulo: o
Volume~0, que trata da notação e do cálculo que os outros dois pressupõem; o
Volume~I, que estabelece os fundamentos de mecânica do contínuo, termofluidos e
sistemas de medição; e o Volume~II, que particulariza tudo isso no compressor
de parafuso isento de óleo de uma planta offshore e na metodologia de Gêmeo
Digital que o acompanha.

A separação em volumes não é administrativa: ela responde a públicos com
pontos de partida diferentes. Quem vem da operação e quer o diagnóstico de uma
máquina real começa no Volume~II e recua sempre que um símbolo o detiver. Quem
vem da formação recente pode abrir no Volume~I. Quem está há anos fora da sala
de aula, ou nunca viu o formalismo, encontra no Volume~0 o percurso curto que
torna os outros dois legíveis. O volume único existe para quem prefere ter o
caminho inteiro à mão, sem trocar de arquivo no meio de uma consulta.

Três decisões de construção merecem registro, porque afetam o modo de usar
este volume. A primeira: a numeração de capítulos, equações, figuras e tabelas
de cada volume é exatamente a da edição avulsa. A Equação~5.5 do Volume~II
continua sendo a Equação~5.5 aqui dentro, e toda citação externa a ela ---
inclusive as que estão no código dos simuladores que acompanham a obra ---
permanece válida. A segunda: os volumes avulsos continuam sendo os documentos
canônicos; esta edição é \emph{gerada} a partir deles, de modo que nunca
divirja do que eles dizem. A terceira: a bibliografia dos três foi unificada
numa lista única ao final, sem duplicatas.

O que a obra não é também merece ser dito. Não é um manual de fabricante nem
um substituto das normas que cita; não traz dados de uma máquina real, e os
números dos estudos de caso são ilustrativos, escolhidos para exercitar as
equações dentro de faixas fisicamente plausíveis. Onde um parâmetro não tem
fonte fechada, o texto diz que não tem, em vez de sugerir uma precisão que não
existe. Essa honestidade sobre o que é medido, o que é inferido e o que é
assumido é, no fim, a condição para que um Gêmeo Digital sirva a uma decisão de
manutenção.

\vspace{1em}
\hfill \textit{O Autor}
